\section{Introduction}\label{intro}
Bid rigging in public procurement is a pervasive and costly problem that undermines competition and leads to artificially inflated prices, ultimately hurting taxpayers and diminishing the quality of public services \citep{porter1993detecting, pesendorfer2000study}. Across industries and countries, numerous cartels have been exposed—ranging from school milk contracts in the United States \citep{porter1993detecting} to large-scale infrastructure projects in Europe and Latin America—which collectively demonstrate the persistent nature of collusive agreements over time. When firms conspire to prearrange who will win a given tender and at what price, the resulting lack of competition raises the costs of government projects, reduces incentives for innovation, and undermines public trust in the integrity of procurement processes. Beyond these immediate economic impacts, collusion can also deter honest businesses from entering the market, thus entrenching cartels and perpetuating a vicious cycle of reduced competition.

Despite stronger antitrust laws and growing sophistication in enforcement, detecting bid rigging remains a significant challenge because explicit collusion is often hidden behind layers of indirect signals and subtle coordination \citep{bajari2003detecting}. Conventional anti-cartel measures—such as investigating suspiciously high bids or relying on whistleblowers—can be reactive and incomplete, because colluding firms endeavor to hide obvious evidence of conspiracy. In the last two decades, researchers have increasingly turned to empirical indicators in auction data to unveil bid coordination. Early work by \citet{porter1993detecting}, for instance, revealed that some construction firms systematically submitted non-competitive bids in New York highway projects, demonstrating how inconsistent pricing patterns can flag collusive behavior. More recent studies apply advanced econometric techniques or big data analytics to procurement datasets, trying to isolate signals that are inconsistent with genuine competitive bidding \citep{conley2016detecting, imhof2018screening}.

In this paper, we employ the methodology proposed by \cite{kawai2023using} and utilize a Regression Discontinuity design to identify patterns indicative of anti-competitive behavior in São Paulo's public procurement data, covering 3.8 million contracts from 2009 to 2019. Our analysis specifically targets auctions decided by narrow margins to detect unusual incumbency advantages and strategic bid timing, which may suggest the presence of anti-competitive schemes among competing firms.

The intuition behind the method hinges on the fact that when two firms bid nearly identical prices, the winner’s edge should be effectively random if all bidders are competing fairly. In other words, under true competition, there is no systematic reason for a narrowly winning firm to differ in characteristics, such as workload, cost structure, frequency of recent wins, from a narrowly losing firm. Hence, if we do observe abrupt differences in these attributes exactly at the point where bids are nearly tied, we can infer that external factors, possibly collusion or corrupt favoritism, are shaping the outcome. While other screening tools often rely on cross-auction patterns or price dispersion, this close-margin design zeroes in on the final, decisive cutoff where a contract is awarded.

In this paper, we apply the close-margin identification strategy to an extensive dataset of public procurement auctions in Brazil between 2009 and 2019. Brazil represents a particularly informative setting for studying collusion and favoritism due to its large-scale procurement programs, varied market structures, and well-documented concerns over corruption in certain sectors \citep{oecd2015, bosio2022public}. Over this period, electronic platforms became the norm for purchasing common goods and services, creating detailed digital records of bids and contract awards. These records enable us to trace whether the same group of firms repeatedly wins contracts by tiny margins or exhibits other patterns suggestive of non-competitive behavior.

Our empirical analysis provides strong evidence of significant discontinuities at the threshold of closely contested auctions, indicative of potential anti-competitive behaviors within public procurement markets. Specifically, we identify two distinct effects: first, firms narrowly winning auctions are approximately 1.8 percentage points more likely to be incumbents—firms that previously secured similar contracts—compared to narrowly losing bidders. This pattern suggests incumbents may gain disproportionate advantages, potentially through strategic bidding or coordinated behaviors. Second, narrowly winning firms demonstrate a substantially higher likelihood, around 6.4 percentage points, of submitting the auction's final bid. This indicates possible strategic timing of bids, hinting at practices consistent with bid leakage or collusive behavior aimed at suppressing genuine competition.
Importantly, these findings are robust after accounting for year and market fixed effects, strengthening the interpretation that observed discontinuities are reflective of deliberate, anti-competitive strategies rather than random occurrences.

The structure of the paper is as follows. Section~\ref{hist} provides an overview of Brazil’s institutional background, describing how public procurement operates under a decentralized yet increasingly digitized framework. Section~\ref{data} describes the dataset and explains our construction of key variables, including the precise definition of the margin between winner and runner-up bids. We then lay out the RDD methodology in Section~\ref{coll}, detailing how we isolate borderline auctions and measure any discontinuities in firm attributes around the winning threshold. Section~\ref{res} presents our main empirical results, showing stark evidence of rotation-type collusion. Section~\ref{res_ar} explores additional analyses and robustness checks, such as splitting the sample by contract size and testing for local manipulation of the margin. Finally, Section~\ref{concl} offers concluding remarks on the policy implications of our findings and potential directions for future research on collusion detection in procurement auctions. 
